\section{Introduction}


E-commerce generates events at high velocity: page views, product interactions, cart modifications, and transactions flow continuously. Extracting actionable insights from this stream requires analytics systems that balance three competing requirements:

\begin{enumerate}
\item \textbf{Freshness}: Insights must reflect recent events. Stale data delays response to problems (site errors, fraud spikes) and opportunities (trending products, effective campaigns).
\item \textbf{Flexibility}: Merchants need ad-hoc queries, not just pre-defined reports. Investigating anomalies requires drilling into arbitrary dimensions.
\item \textbf{Scale}: Event volumes grow with business success. Systems must scale horizontally without architectural changes.
\end{enumerate}

Traditional approaches force trade-offs. Data warehouses (Redshift, BigQuery) provide flexibility but with batch latency. Stream processors (Storm, Flink) offer freshness but limited query expressiveness. Lambda architectures \cite{marz2015} combine both but double operational complexity.

Verus achieves all three requirements through a unified architecture. Events stream into a real-time aggregation layer that maintains pre-computed rollups for common queries. Simultaneously, events land in columnar storage for ad-hoc analysis. A query planner routes requests to the optimal layer based on query characteristics.

Our contributions are:
\begin{enumerate}
\item A hybrid architecture combining stream processing and columnar storage with automatic query routing.
\item Materialized view maintenance for sub-second dashboard queries.
\item Cohort analysis operators for retention and lifetime value computation.
\item Anomaly detection integrated into the analytics pipeline.
\end{enumerate}

